\documentclass[11pt,a4paper,openany]{ctexbook}

\usepackage[a4paper,left=22mm,right=22mm,top=20mm,bottom=22mm]{geometry}
\usepackage{amsmath,amssymb}
\usepackage{booktabs,longtable,array}
\usepackage{enumitem}
\usepackage{xcolor}
\usepackage{hyperref}
\usepackage{listings}
\usepackage{fancyhdr}
\usepackage{graphicx}

\hypersetup{
  colorlinks=true,
  linkcolor=blue!55!black,
  urlcolor=blue!55!black,
  citecolor=blue!55!black,
  pdftitle={CCUSDT 工程回测教材 v0.2},
  pdfauthor={Jiang / Codex}
}

\setlength{\parindent}{2em}
\setlength{\parskip}{0.46em}
\setlength{\headheight}{14pt}
\linespread{1.24}
\raggedbottom

\pagestyle{fancy}
\fancyhf{}
\fancyhead[L]{CCUSDT 工程回测教材}
\fancyhead[R]{简单状态机，从数据到自己回测}
\fancyfoot[C]{\thepage}

\definecolor{codebg}{RGB}{247,248,247}
\definecolor{rulegray}{RGB}{210,215,210}
\definecolor{deepgreen}{RGB}{18,88,62}
\definecolor{softblue}{RGB}{36,84,128}
\definecolor{warnred}{RGB}{130,45,45}

\lstdefinestyle{code}{
  basicstyle=\ttfamily\small,
  backgroundcolor=\color{codebg},
  frame=single,
  rulecolor=\color{rulegray},
  breaklines=true,
  columns=fullflexible,
  keepspaces=true,
  showstringspaces=false,
  tabsize=2
}
\lstset{style=code}

\newcommand{\code}[1]{\texttt{\detokenize{#1}}}
\urlstyle{tt}
\newcommand{\file}[1]{\nolinkurl{#1}}
\newcommand{\term}[1]{\textbf{#1}}
\newcommand{\truth}[1]{\par\noindent\textcolor{deepgreen}{\textbf{真相：}}#1\par}
\newcommand{\warningline}[1]{\par\noindent\textcolor{warnred}{\textbf{边界：}}#1\par}
\newcommand{\checkpoint}[1]{\par\noindent\textcolor{softblue}{\textbf{检查点：}}#1\par}
\newcommand{\ds}{\displaystyle}

\title{\Huge CCUSDT 工程回测教材\\[6pt]
\Large L2 第一性原理到完整 fast/Runner 系统}
\author{面向 Python 会一点、Rust 新人的未来自己与新 Codex}
\date{v0.2 / 2026-06-03}

\begin{document}
\frontmatter
\maketitle

\begin{center}
\begin{minipage}{0.84\textwidth}
本书不是 API 调用大全，也不是把简单策略包装成复杂架构的系统白皮书。
它只解决一个朴素问题：
\begin{center}
\term{怎样从 L2 订单簿和成交开始，完整掌握当前这个简单策略系统？}
\end{center}
当前策略本身不复杂。第一章先用 L2 订单簿把策略定义全部讲清楚；第二章开始再讲
数据如何转换成 \code{decision_frame}，以及 Python fast、Python Bot、Rust Runner
分别怎样执行同一个状态机。
\end{minipage}
\end{center}

\newpage
\tableofcontents
\newpage

\mainmatter
\setcounter{chapter}{-1}

\input{docs/books/ccusdt-engineering-textbook/chapters/v02-00-remove-mystery}
\input{docs/books/ccusdt-engineering-textbook/chapters/v02-01-l2-strategy}
\input{docs/books/ccusdt-engineering-textbook/chapters/v02-02-data-shapes}
\input{docs/books/ccusdt-engineering-textbook/chapters/v02-03-decision-frame-pipeline}
\input{docs/books/ccusdt-engineering-textbook/chapters/v02-04-fast-backtest}
\input{docs/books/ccusdt-engineering-textbook/chapters/v02-05-artifacts}
\input{docs/books/ccusdt-engineering-textbook/chapters/v02-06-python-bot}
\input{docs/books/ccusdt-engineering-textbook/chapters/v02-07-rust-runner}
\input{docs/books/ccusdt-engineering-textbook/chapters/v02-08-strict-order-life}
\input{docs/books/ccusdt-engineering-textbook/chapters/v02-09-fast-vs-runner}
\input{docs/books/ccusdt-engineering-textbook/chapters/v02-10-one-day-backtest}
\input{docs/books/ccusdt-engineering-textbook/chapters/v02-11-safe-experiment}
\input{docs/books/ccusdt-engineering-textbook/chapters/v02-12-not-implemented}
\appendix
\input{docs/books/ccusdt-engineering-textbook/chapters/v02-appendix-fields}
\input{docs/books/ccusdt-engineering-textbook/chapters/v02-appendix-code}
\input{docs/books/ccusdt-engineering-textbook/chapters/v02-appendix-commands}

\end{document}
